\subsection{Protocolo de test de usabilidad}

El protocolo propone evaluar si los recorridos de Carer-Record permiten al personal asistencial consultar contexto autorizado, registrar notas propias y revisar la asistencia y los resúmenes de período. La evaluación emplea datos sintéticos y no se afirma como ejecutada.

\subsubsection*{Objetivos del test} ~

El test busca observar la completitud de las tareas, el tiempo empleado, los errores y la comprensión de las vistas. También recoge la percepción general de usabilidad para orientar iteraciones de diseño. Los indicadores y umbrales descritos son criterios propuestos para esta evaluación, no resultados ni requisitos funcionales.

\subsubsection*{Perfil de participantes} ~

Se propone invitar a un mínimo de cinco participantes que correspondan al perfil de personal asistencial: médico tratante, médico de turno o profesional de enfermería con al menos un año de experiencia hospitalaria y uso habitual de smartphone. La participación es voluntaria y requiere consentimiento informado antes de comenzar. Este perfil no implica que se hayan reclutado o evaluado participantes.

\subsubsection*{Tareas a evaluar} ~

El Cuadro~\ref{tab:tareas_test_usabilidad} define seis tareas para recorrer el acceso, la consulta, el registro y la revisión de asistencia. Los tiempos son objetivos orientativos para observar el recorrido; no son garantías de rendimiento ni condiciones de aceptación del producto.

\begin{table}[H]
\caption{Tareas propuestas para el test de usabilidad}
\label{tab:tareas_test_usabilidad}
\begin{tabularx}{\textwidth}{|l|X|X|c|}
\hline
\textbf{N} & \textbf{Tarea} & \textbf{Criterio observable} & \textbf{Tiempo objetivo} \\
\hline
T-01 & Iniciar sesión y abrir una hospitalización sintética autorizada. & Se muestra el contexto permitido; un acceso sin autorización es denegado. & 30 s \\
\hline
T-02 & Consultar información de la FCE simulada y notas propias en la bitácora. & Se identifican las procedencias y el orden temporal de la información. & 2 min \\
\hline
T-03 & Redactar y confirmar una nota mediante texto. & La nota revisada queda asociada a la persona autora y a la hospitalización. & 60 s \\
\hline
T-04 & Dictar una nota, revisar la transcripción local y corregirla antes de confirmar. & Solo el texto revisado y confirmado queda registrado. & 60 s \\
\hline
T-05 & Solicitar asistencia de redacción y aceptar, editar o descartar la propuesta. & La decisión corresponde a la persona y el texto de origen no se sustituye automáticamente. & 90 s \\
\hline
T-06 & Seleccionar un período con varios eventos y notas y revisar su resumen. & El resumen identifica el período y sus fuentes, y no se publica como nota ni se escribe en la FCE. & 45 s \\
\hline
\end{tabularx}
\end{table}

\subsubsection*{Métricas y criterios de iteración} ~

Para cada tarea se propone registrar:

\begin{itemize}
    \item \textit{Completitud:} proporción de participantes que completa la tarea sin ayuda. Una completitud inferior al 80\% se propone como señal para revisar el diseño de esa tarea.
    \item \textit{Tiempo:} duración observada comparada con el objetivo orientativo del cuadro. Un promedio que supere el objetivo en más de 50\% se propone como señal para revisar el recorrido, no como incumplimiento del producto.
    \item \textit{Errores:} tipo y cantidad, distinguiendo errores recuperables de errores bloqueantes que requieren ayuda o impiden terminar.
    \item \textit{Usabilidad percibida:} aplicación del cuestionario System Usability Scale (SUS) de diez ítems al terminar; un resultado inferior a 70 puntos se propone como señal para iterar la interfaz~\cite{ref_brooke_sus}.
\end{itemize}

Los umbrales orientan decisiones de diseño y deben interpretarse junto con las observaciones y condiciones de aplicación; no acreditan por sí mismos la aceptación del producto. Una nueva ronda puede evaluar los cambios introducidos. Esta evaluación de usabilidad no sustituye la verificación técnica ni una revisión clínica diferenciada del contenido generado.

\subsubsection*{Protocolo de aplicación} ~

La prueba se propone en formato presencial, con una persona moderadora. Tras obtener el consentimiento, se entrega un escenario sintético y se permite que cada participante complete las tareas sin instrucciones sobre cómo usar la interfaz. La persona moderadora observa y registra pasos, tiempos y errores, e interviene solo si existe bloqueo o si la persona solicita detener la prueba. No se requiere grabar audio o video; cualquier registro adicional exige consentimiento explícito.

El escenario presenta a una persona asistencial que necesita consultar una hospitalización autorizada, revisar la FCE simulada y las notas propias, registrar contexto mediante texto y dictado, revisar una propuesta de redacción y generar un resumen de un período seleccionado. El resumen se revisa en Carer-Record y no se copia ni se escribe en la ficha oficial. El intervalo seleccionado no representa un turno.

Al terminar, se propone una entrevista breve sobre los pasos más confusos, los cambios sugeridos, la diferencia entre la FCE y las notas propias y si la persona usaría el sistema en su trabajo. La ejecución y sus resultados se documentarán si la evaluación se realiza; este protocolo no informa resultados.
